\chapter{Cadre conceptuel et étude de l'existant}
\label{ch:cadre-conceptuel}
\minitoc
\clearpage

\section{Introduction}

Avant d'aborder les choix techniques, il est nécessaire de replacer PharmacieConnect dans le contexte
opérationnel qui a motivé sa réalisation. Ce chapitre part du fonctionnement des gardes pharmaceutiques,
présente le cadre du stage externe, puis examine les difficultés rencontrées lorsqu'une information de
garde doit être cherchée, corrigée ou publiée rapidement. Il introduit enfin la réponse retenue et la
manière dont le projet a été organisé.

\section{Présentation du domaine pharmaceutique}

Le domaine pharmaceutique ne se limite pas à la vente de médicaments : il participe aussi à la continuité
des soins, notamment lorsque les pharmacies de garde assurent une disponibilité en dehors des horaires
ordinaires \cite{cnopt_tableau_garde,cnopt_cropt_gardes}. Dans le cadre de PharmacieConnect, cette réalité métier se traduit par des
données précises à gérer : pharmacie, adresse, téléphone, gouvernorat, position GPS, date de garde, horaire
et type de garde.

La valeur de ces données dépend directement de leur actualité. Une pharmacie mal localisée, une date de
garde incorrecte ou un numéro de téléphone obsolète peut empêcher l'utilisateur de trouver rapidement une
solution. C'est pourquoi le projet traite la maintenance des données comme une exigence centrale : les
imports CSV sont validés, les comptes administratifs sont contrôlés par rôle et les actions sensibles sont
enregistrées dans les journaux d'audit.

\section{Contexte général}

PharmacieConnect regroupe les informations liées aux pharmacies, aux plannings de garde, aux médicaments et
aux premiers secours. La plateforme distingue deux usages complémentaires. Le citoyen utilise l'application
mobile pour chercher une pharmacie, consulter les gardes, accéder au catalogue de médicaments ou poser une
question à l'assistant. L'administrateur et les utilisateurs autorisés utilisent le tableau de bord web pour
importer, corriger, contrôler et suivre les données diffusées.

Le backend constitue la couche commune entre ces interfaces. Il expose les endpoints publics nécessaires à
la consultation mobile et les endpoints protégés utilisés par l'administration. Cette séparation évite que
les règles de validation, de rôle et de journalisation soient dupliquées dans les clients.

\section{Cadre du stage externe et organisme d'accueil}

Le projet a été réalisé dans le cadre d'un stage externe au sein de \textbf{IT-STORM Consulting}. Le sujet
officiel concerne le développement d'une application mobile permettant aux citoyens de consulter les
pharmacies de garde en Tunisie. Les principaux besoins portent sur la recherche par géolocalisation ou
par saisie manuelle, l'affichage en liste et sur carte, la consultation des coordonnées et l'accès rapide
aux informations utiles.

Le stage s'est déroulé du 07/02/2026 au 07/06/2026, sous l'encadrement professionnel de
\textbf{Achraf BCHIRI}. Ce cadre externe a orienté le projet vers une solution exploitable, structurée
autour d'une API REST~\cite{fielding}, d'une base de données, d'une application mobile et d'une interface
d'administration pour la mise à jour des données.

\begin{table}[H]
    \centering
    \caption{Cadre institutionnel du projet\\ \small Source : Auteur}
    \begin{tabular}{L{4.5cm}L{8.5cm}}
        \toprule
        \textbf{Élément} & \textbf{Description} \\
        \midrule
        Organisme d'accueil & IT-STORM Consulting. \\
        Adresse & 12 Rue Jean Pacilly, 91120 Palaiseau, France. \\
        Encadrant professionnel & Achraf BCHIRI. \\
        Période & Du 07/02/2026 au 07/06/2026. \\
        Nature du projet & Projet de fin d'études externe orienté conception et réalisation logicielle. \\
        Domaine & Santé numérique, accès aux pharmacies, gardes et données pharmaceutiques. \\
        Périmètre technique & API backend, tableau de bord administrateur, application mobile et base de données. \\
        Encadrement & Suivi professionnel, suivi académique, revues d'avancement et validation des livrables. \\
        \bottomrule
    \end{tabular}
\end{table}

Ce cadre a orienté la démarche vers un prototype complet et vérifiable. La priorité a été donnée aux
fonctionnalités essentielles de recherche mobile, à la séparation des responsabilités, à la fiabilité des
données et à la traçabilité des actions administratives.

\section{Problèmes et défis}

Les problèmes observés concernent principalement la disponibilité, la fiabilité et l'exploitation des
données. Les informations peuvent être publiées sous des formats différents, mises à jour avec retard ou
présentées d'une manière peu adaptée à une recherche mobile.

Les défis retenus pour le projet sont les suivants :

\begin{itemize}
    \item centraliser les informations sur les pharmacies et les gardes ;
    \item rendre la recherche simple pour un utilisateur mobile ;
    \item permettre la consultation par date et par proximité ;
    \item faciliter l'import de volumes de données à travers des fichiers \gls{csv} ;
    \item sécuriser l'administration et limiter les actions selon les rôles ;
    \item garder une trace des actions sensibles pour améliorer la traçabilité.
\end{itemize}

\section{Étude de l'existant}

L'étude de l'existant a été orientée par les besoins du projet. Il ne s'agit pas seulement de vérifier si
une solution affiche une pharmacie, mais d'évaluer sa capacité à répondre à plusieurs critères : disponibilité
des gardes, usage mobile, géolocalisation, fraîcheur des données, workflow de mise à jour, administration,
authentification, import de données, auditabilité et couverture tunisienne.

Les listes statiques diffusées sous forme de page web, d'image ou de document restent faciles à partager,
mais elles deviennent limitées lorsqu'un utilisateur veut filtrer par gouvernorat, choisir une date précise
ou ouvrir l'adresse sur une carte. Elles ne renseignent pas toujours le processus de mise à jour ni l'origine
de la correction lorsqu'une donnée change \cite{cnopt_tableau_garde,medtn}.

Les sites spécialisés et les applications mobiles offrent une expérience plus adaptée à la consultation.
Cependant, leur efficacité dépend de la qualité de la base de données, de la fréquence de mise à jour et de
la présence d'un espace d'administration capable de contrôler les imports, les droits et les corrections.
Cette distinction est importante pour PharmacieConnect, car le projet vise à couvrir à la fois la
consultation citoyenne et la gestion opérationnelle des données.

\section{Analyse du marché et solutions existantes}

Pour situer PharmacieConnect, plusieurs familles de solutions ont été comparées. Le Conseil National de
l'Ordre des Pharmaciens de Tunisie constitue une source institutionnelle liée à l'organisation de la
profession \cite{cnopt}. Des annuaires médicaux comme Med.tn publient des informations pratiques sur les
pharmacies de garde \cite{medtn}. Les solutions cartographiques, telles que Google Maps et OpenStreetMap,
facilitent la localisation d'une pharmacie, mais elles ne constituent pas un outil métier de planification
des gardes ni de validation des imports \cite{googlemaps,openstreetmap}. Enfin, l'application Apothicare,
associée au Conseil national de l'ordre des pharmaciens de Tunisie, illustre l'intérêt d'un accès mobile aux
pharmacies ouvertes ou de garde \cite{apothicare}.

La comparaison suivante reste fonctionnelle : elle évalue les capacités observables ou annoncées par les
solutions citées. Lorsque l'information n'est pas confirmée par une source publique, la mention doit être
vérifiée avant la version finale \cite{cnopt,medtn,googlemaps,openstreetmap,apothicare}.

\begin{table}[H]
    \centering
    \caption{Comparaison fonctionnelle des solutions existantes\\ \small Source : Auteur, à partir de \cite{cnopt,medtn,googlemaps,openstreetmap,apothicare}}
    \label{tab:comparaison-solutions}
    \scriptsize
    \begin{tabular}{L{2.6cm}C{1.1cm}C{1.3cm}C{1.5cm}C{1.6cm}C{1.4cm}C{1.7cm}}
        \toprule
        \textbf{Solution} & \textbf{API REST} & \textbf{App mobile} & \textbf{Gestion planning} & \textbf{Auth.} & \textbf{Import données} & \textbf{Couverture} \\
        \midrule
        Sources officielles (CNOPT) & Non & Non & Partiel & Non & Non & Oui \\
        Annuaires médicaux (Med.tn) & Non & Non & Partiel & Non & Non & Oui \\
        Cartographie (Google Maps / OSM) & Oui & Oui & Non & Non & Non & Partiel \\
        Apps mobiles (Apothicare) & Non & Oui & Oui & Partiel & Non & Oui \\
        Listes statiques & Non & Non & Non & Non & Non & Partiel \\
        \shortstack[l]{\textbf{Pharmacie}\\\textbf{Connect}} & \textbf{Oui} & \textbf{Oui} & \textbf{Oui} & \textbf{Oui} & \textbf{Oui} & \textbf{Oui} \\
        \bottomrule
    \end{tabular}
\end{table}

\section{Critique de l'existant}

L'étude de l'existant montre que le principal risque ne concerne pas uniquement l'affichage des pharmacies.
Il concerne aussi la maîtrise de la donnée. Une solution peut permettre la consultation, mais rester limitée
si elle ne précise pas comment les plannings sont importés, corrigés, validés et tracés. Pour un utilisateur
mobile, cette limite se traduit par une recherche moins fiable. Pour l'administrateur, elle se traduit par
une difficulté à savoir qui a modifié une donnée, quand et à partir de quelle source.

Cette difficulté augmente avec le volume de données et le nombre d'intervenants. Une saisie manuelle devient
insuffisante lorsque plusieurs assistants régionaux doivent gérer des pharmacies, des gardes et des fichiers
CSV. Sans validation des colonnes, contrôle des doublons, restriction des rôles et journalisation, la base
peut devenir incohérente sans que l'origine de l'erreur soit identifiable.

\section{Solution proposée : PharmacieConnect}

\textbf{PharmacieConnect} est proposée comme une réponse intégrée aux limites observées. La solution est
structurée autour de trois composants principaux, auxquels s'ajoute un assistant conversationnel de premiers
secours intégré au parcours mobile :

\begin{itemize}
    \item une API backend FastAPI qui centralise les données, gère l'authentification, applique les règles
    métier et expose les endpoints \gls{rest} ;
    \item un tableau de bord administrateur React/Vite qui offre aux administrateurs un espace de gestion
    pour les pharmacies, les gardes, les médicaments, les imports \gls{csv}, les assistants régionaux, les
    indicateurs et les journaux d'audit ;
    \item une application mobile Expo/React Native centrée sur le citoyen, avec la recherche de
    pharmacies, la consultation des gardes, la géolocalisation, l'accès au catalogue de médicaments et
    l'assistant de premiers secours.
\end{itemize}

Cette organisation répond aux limites relevées dans l'existant : dispersion des informations, difficulté de
mise à jour, absence de traçabilité et manque de séparation entre consultation publique et gestion
administrative. Les données sont centralisées dans le backend, les accès sont contrôlés selon les rôles et
les interfaces sont adaptées aux deux usages principaux : la consultation mobile et l'administration. L'assistant
ajoute un usage complémentaire en orientant l'utilisateur vers des réponses prudentes de premiers secours,
sans remplacer un professionnel de santé ni les services d'urgence.

\section{Objectifs}

Les objectifs fonctionnels et techniques du projet sont :

\begin{itemize}
    \item offrir une recherche rapide des pharmacies ;
    \item afficher les pharmacies de garde selon une date donnée ;
    \item permettre une recherche par proximité géographique ;
    \item administrer les pharmacies, les gardes, les médicaments et les comptes autorisés ;
    \item importer des fichiers \gls{csv} en contrôlant la qualité des données ;
    \item sécuriser les accès à travers l'authentification, les rôles et la validation côté serveur ;
    \item produire des journaux d'audit et des indicateurs utiles à l'administration.
\end{itemize}

\section{Méthodologie de développement}

Le projet a adopté une organisation par sprints inspirée de Scrum \cite{scrumguide}, adaptée au contexte
d'un PFE. Chaque sprint a été associé à un objectif concret et à des livrables vérifiables. Les arbitrages
ont porté sur les modules réellement développés : contrats d'API, imports CSV, authentification, gestion des
rôles, pages d'administration, écrans mobiles et tests de validation.

\begin{table}[H]
    \centering
    \caption{Organisation des sprints de réalisation\\ \small Source : Auteur}
    \label{tab:sprints-scrum}
    \begin{tabular}{L{2.6cm}L{5.1cm}L{5.3cm}}
        \toprule
        \textbf{Sprint} & \textbf{Objectif} & \textbf{Livrables principaux} \\
        \midrule
        Sprint 1 & Analyse et cadrage & Besoins, acteurs, règles de gestion, étude de l'existant. \\
        Sprint 2 & Conception technique & Architecture, modèle de données, sécurité, contrats API. \\
        Sprint 3 & Backend & Routes FastAPI, services métier, authentification, imports et audit. \\
        Sprint 4 & Administration web & Pages React/Vite, tableaux, formulaires, indicateurs et gestion des droits. \\
        Sprint 5 & Application mobile & Recherche, carte, gardes, médicaments, paramètres et langues. \\
        Sprint 6 & Intégration et validation & Tests, build, corrections, rapport et annexes. \\
        \bottomrule
    \end{tabular}
\end{table}

Cette méthode a limité les risques liés à une intégration tardive. Les modules sensibles, notamment
l'authentification, les imports CSV, les restrictions régionales, les journaux d'audit et la recherche mobile,
ont été vérifiés progressivement par des tests automatisés ou par des scénarios fonctionnels.

\section{Outils de gestion de projet}
\label{sec:outils-gestion}

La conduite du projet s'est appuyée sur des outils de suivi et de versionnement afin de relier les besoins,
les tâches et le code livré. \textbf{Jira} a été utilisé pour organiser les tickets, les sprints et les
anomalies \cite{atlassian_jira_sprints}. \textbf{Bitbucket} a été utilisé pour héberger le dépôt Git, suivre les
branches et préparer les fusions de code \cite{atlassian_bitbucket_guides}. Cette organisation a fourni une traçabilité
entre les fonctionnalités prévues et les modules effectivement réalisés.

\subsection{Jira pour la planification et le suivi}

Jira a été utilisé pour la planification des sprints et le suivi des tâches. Chaque sprint défini dans le
tableau~\ref{tab:sprints-scrum} a été décliné en tickets, puis en sous-tâches lorsque le module le
nécessitait. Les tickets ont notamment couvert les imports CSV, les pages d'administration, les restrictions
régionales, les tests et les corrections d'interface.

\begin{itemize}
    \item un tableau Kanban/Scrum organisé en quatre colonnes : \emph{À faire}, \emph{En cours},
    \emph{En révision}, \emph{Terminé} ;
    \item un suivi des anomalies (\emph{bugs}) et des demandes d'évolution sous forme de tickets Jira,
    avec attribution d'un niveau de priorité ;
    \item l'utilisation du \emph{burndown chart} et de la vélocité de sprint pour mesurer
    l'avancement à chaque cycle.
\end{itemize}

\subsection{Bitbucket pour le versionnement du code}

Bitbucket a hébergé le dépôt Git du projet. L'organisation du dépôt a suivi une logique inspirée de Gitflow
\cite{atlassian_gitflow} :

\begin{itemize}
    \item branche \texttt{main} pour les versions stables ;
    \item branche \texttt{develop} pour l'intégration courante ;
    \item branches \texttt{feature/*} pour les nouvelles fonctionnalités ;
    \item branches \texttt{bugfix/*} pour les corrections d'anomalies.
\end{itemize}

Chaque branche de fonctionnalité correspondait à un ticket Jira, en respectant la convention de nommage
\texttt{feature/PC-XX-description}. Les fusions vers \texttt{develop} étaient soumises à une
\emph{pull request} faisant l'objet d'une revue de code. Les messages de \emph{commit} référençaient
systématiquement le numéro de ticket Jira (par exemple \texttt{PC-42: ajout import CSV gardes}), ce qui
a garanti une traçabilité de bout en bout. Le dépôt monorepo regroupait quatre dossiers principaux :
\texttt{backend/}, \texttt{admin/}, \texttt{mobile/} et \texttt{chatbot/}.

\begin{table}[H]
    \centering
    \caption{Outils de gestion de projet utilisés\\ \small Source : Auteur}
    \label{tab:outils-gestion}
    \begin{tabular}{L{2.4cm}L{4.6cm}L{6.0cm}}
        \toprule
        \textbf{Outil} & \textbf{Rôle} & \textbf{Utilisation principale} \\
        \midrule
        Jira & Gestion des sprints et des tâches & Planification Scrum, suivi des user stories, gestion des bugs et des évolutions, tableau Kanban et burndown chart. \\
        Bitbucket & Versionnement du code source & Hébergement Git, branches Gitflow, pull requests, revue de code et traçabilité par référence aux tickets Jira. \\
        \bottomrule
    \end{tabular}
\end{table}

\subsection{Intégration Jira--Bitbucket}

L'intégration entre Jira et Bitbucket a permis de relier les tickets aux branches, aux commits et aux
\emph{pull requests} \cite{atlassian_jira_bitbucket}. Dans le cadre du projet, cette liaison a facilité le suivi entre
une exigence fonctionnelle, par exemple l'import des gardes ou la consultation des audit logs, et les
modifications de code correspondantes.

\section{Conclusion}

Ce premier chapitre a permis de préciser le contexte du projet, le domaine pharmaceutique et les limites
des solutions existantes. Il montre également pourquoi PharmacieConnect repose sur une séparation entre
backend, administration web et application mobile. Le chapitre suivant détaille les besoins qui découlent
de ce cadrage.
